%% Methodology section for the SCMK anomaly-detection paper.
%% Built on the Elsevier elsarticle class (numbered citation style).
%% Compile with: pdflatex method.tex  (run twice for cross-references).

\documentclass[preprint,12pt]{elsarticle}

\usepackage{amssymb}
\usepackage{amsmath}
\usepackage{amsthm}
\usepackage{algorithm}
\usepackage{algpseudocode}
\usepackage{bm}

%% Lightweight theorem environments for the scatter-equivalence result.
\newtheorem{proposition}{Proposition}
\newtheorem{remark}{Remark}

%% Convenience macros.
\newcommand{\R}{\mathbb{R}}
\newcommand{\Sph}{\mathbb{S}}
\newcommand{\norm}[1]{\left\lVert #1 \right\rVert}
\newcommand{\inner}[2]{\left\langle #1,\, #2 \right\rangle}
\newcommand{\Xnorm}{\mathcal{X}_{\mathrm{n}}}
\newcommand{\Lcross}{\mathcal{L}_{\mathrm{con}}}
\newcommand{\Lscat}{\mathcal{L}_{\mathrm{scat}}}
\newcommand{\Ltot}{\mathcal{L}}
\newcommand{\hh}{\bm{h}}
\newcommand{\zz}{\bm{z}}
\newcommand{\xx}{\bm{x}}
\newcommand{\ww}{\bm{w}}
\newcommand{\WW}{\bm{W}}
\newcommand{\mmu}{\bm{\mu}}

\journal{Pattern Recognition}

\begin{document}

\begin{frontmatter}

\title{Scatter-Regularized Cross-Multi-Kernel Contrastive
Learning for One-Class Anomaly Detection}

\author{}
\affiliation{organization={},
            addressline={},
            city={},
            postcode={},
            state={},
            country={}}

\begin{abstract}
We present the methodology of a scatter-regularized cross-multi-kernel
contrastive framework (SCMK) for one-class anomaly detection. The full
manuscript is in preparation; this file contains the Methodology section.
\end{abstract}

\begin{keyword}
anomaly detection \sep multiple kernel learning \sep contrastive learning
\sep one-class classification \sep representation learning
\end{keyword}

\end{frontmatter}

%% ======================================================================
\section{Methodology}
\label{sec:method}

\subsection{Problem formulation and notation}
\label{subsec:problem}

We address \emph{one-class} (semi-supervised) anomaly detection. Let
$\mathcal{X}=\{\xx_i\}_{i=1}^{N}$ with $\xx_i\in\R^{D}$ denote the data after
preprocessing,\footnote{Nominal attributes are one-hot encoded and numerical
attributes are min--max scaled to $[0,1]$, so that all feature dimensions share
a comparable range.} and let $y_i\in\{0,1\}$ be the latent label, where $y_i=0$
marks a normal instance and $y_i=1$ an anomaly. Only the normal subset
$\Xnorm=\{\xx_i : y_i=0\}$, of size $N_{\mathrm{n}}=|\Xnorm|$, is available at
training time; anomalies are neither observed nor simulated. The goal is to
learn a scoring function $a:\R^{D}\!\to\!\R$ such that $a(\xx)$ is large for
anomalies and small for normal points.

Our framework, termed \emph{scatter-regularized cross-multi-kernel}
contrastive learning (SCMK), couples a multiple-kernel contrastive encoder
with a one-class detection head. It comprises four components, summarised in
Algorithm~\ref{alg:scmk}:
(i) a multi-scale kernel bank that injects geometric priors at several
resolutions (Section~\ref{subsec:kernels});
(ii) a cross-kernel contrastive objective that aligns the views produced by
different kernels (Section~\ref{subsec:contrastive});
(iii) a multi-kernel scatter regularizer that explicitly compacts the normal
class in every kernel-induced embedding (Section~\ref{subsec:scatter}); and
(iv) a linear one-class support vector machine (OC-SVM) that scores anomalies
in the concatenated representation (Section~\ref{subsec:detection}).

\subsection{Multi-scale kernel bank}
\label{subsec:kernels}

To endow the encoder with multi-resolution sensitivity, we instantiate a bank
of $K$ Gaussian kernels whose bandwidths are anchored to the intrinsic scale of
the normal data. Let
\begin{equation}
\mathrm{med}
=\operatorname{median}\!\big\{\norm{\xx_i-\xx_j}_2 :
\xx_i,\xx_j\in\Xnorm,\ i\neq j\big\}
\label{eq:median}
\end{equation}
be the median pairwise Euclidean distance estimated \emph{from normal samples
only}, which prevents contamination of the scale statistic by outliers. The
$k$-th kernel is
\begin{equation}
\kappa_k(\bm{a},\bm{b})
=\exp\!\left(-\frac{\norm{\bm{a}-\bm{b}}_2^{2}}{2\,t_k^{2}}\right),
\qquad
t_k = r_k\cdot\mathrm{med},
\label{eq:gaussian}
\end{equation}
with multiplicative ratios $\{r_k\}_{k=1}^{K}=\{0.1,0.5,1.0,2.0,5.0\}$
($K=5$). Small ratios produce sharp kernels that resolve fine-grained local
neighbourhoods, whereas large ratios produce broad kernels that capture the
global manifold geometry. The bank therefore provides $K$ complementary
similarity measures that the contrastive stage will reconcile.

\subsection{Cross-kernel contrastive representation learning}
\label{subsec:contrastive}

\paragraph{Per-kernel embeddings.}
For each kernel we attach an independent, bias-free linear projection head
$g_k(\xx)=\WW_k\xx$ with $\WW_k\in\R^{d\times D}$, followed by
$\ell_2$ normalization,
\begin{equation}
\hh_{k,i}
=\frac{\WW_k\xx_i}{\norm{\WW_k\xx_i}_2}\in\Sph^{d-1},
\qquad k=1,\dots,K,
\label{eq:embedding}
\end{equation}
so that every embedding lives on the unit hypersphere $\Sph^{d-1}$. Removing the
bias keeps the representation centred at the origin, which is consistent with
the origin-separating hyperplane assumption of the downstream OC-SVM, while the
normalization equalises the scale across kernels.

\paragraph{Cross-kernel positive pairs.}
Whereas conventional contrastive learning forms positive pairs through
stochastic data augmentation, we treat the \emph{kernel index as the view}: the
two embeddings $\hh_{k,i}$ and $\hh_{l,i}$ of the \emph{same} sample under
\emph{different} kernels constitute a positive pair, and embeddings of distinct
samples are negatives. Consider a mini-batch of $B$ normal samples and a kernel
pair $(k,l)$. Stacking the two views yields
$\bm{F}^{(kl)}=[\hh_{k,1},\dots,\hh_{k,B},\hh_{l,1},\dots,\hh_{l,B}]^{\!\top}
\in\R^{2B\times d}$,
whose rows are indexed by $m,n\in\{1,\dots,2B\}$. The pairwise affinity is the
\emph{symmetrised kernel response} evaluated on the embeddings,
\begin{equation}
s^{(kl)}_{mn}
=\tfrac{1}{2}\big[\kappa_k(\bm{f}_m,\bm{f}_n)
+\kappa_l(\bm{f}_m,\bm{f}_n)\big].
\label{eq:affinity}
\end{equation}
Because the embeddings are unit-norm, the Gaussian response reduces to a
monotone function of cosine similarity,
$\kappa_k(\bm{f}_m,\bm{f}_n)
=\exp\!\big(-(1-\inner{\bm{f}_m}{\bm{f}_n})/t_k^{2}\big)$,
so $t_k^{2}$ plays the role of a contrastive temperature that is automatically
adapted to the data scale through Eq.~\eqref{eq:median}.

\paragraph{Cross-kernel InfoNCE.}
For anchor $n$, let $\pi(n)$ index its cross-view positive (i.e.\ $\pi(n)=n+B$
for $n\le B$ and $\pi(n)=n-B$ otherwise). The InfoNCE loss for the pair $(k,l)$
is
\begin{equation}
\Lcross^{(kl)}
=-\frac{1}{2B}\sum_{n=1}^{2B}
\log\frac{\exp\!\big(s^{(kl)}_{n,\pi(n)}\big)}
{\displaystyle\sum_{m\neq n}\exp\!\big(s^{(kl)}_{nm}\big)},
\label{eq:infonce}
\end{equation}
where the denominator excludes the self-similarity term $m=n$. Averaging over
all $\binom{K}{2}$ unordered kernel pairs gives the cross-kernel contrastive
objective
\begin{equation}
\Lcross
=\frac{1}{\binom{K}{2}}
\sum_{1\le k<l\le K}\Lcross^{(kl)}.
\label{eq:lcross}
\end{equation}
Minimising Eq.~\eqref{eq:lcross} drives the encoder to a \emph{view-invariant}
representation in which all kernels agree on the embedding of each normal
sample, i.e.\ $\hh_{k,i}\!\approx\!\hh_{l,i}$, while keeping distinct samples
mutually discriminable through the negative terms.

\subsection{Multi-kernel scatter regularization}
\label{subsec:scatter}

\paragraph{Motivation.}
The contrastive objective~\eqref{eq:lcross} enforces \emph{cross-kernel
consistency} but does not constrain the \emph{absolute compactness} of the
normal class: it is invariant to a rotation that spreads normal embeddings
across the sphere as long as the per-sample views remain aligned. A one-class
detector, however, benefits directly from a tightly concentrated normal region.
We therefore introduce a regularizer that minimises the dispersion of the normal
class within every kernel-induced embedding space, drawing on the one-class
specialisation of centroid multiple-kernel $k$-means (CMKKM).

\paragraph{Definition.}
Let $\mmu_k=\frac{1}{B}\sum_{i=1}^{B}\hh_{k,i}$ be the empirical centroid of the
normal embeddings under kernel $k$. The multi-kernel scatter penalty is
\begin{equation}
\Lscat
=-\frac{1}{K}\sum_{k=1}^{K}\norm{\mmu_k}_2^{2}.
\label{eq:lscat}
\end{equation}
It is computed in $O(KBd)$ time --- a single mean per kernel --- and thus adds
negligible overhead.

\paragraph{Exact equivalence to within-class scatter.}
The next result shows that Eq.~\eqref{eq:lscat} is not a heuristic surrogate but
is \emph{exactly} (up to an additive constant) the average within-class scatter
of the normal embeddings.

\begin{proposition}[Scatter equivalence]
\label{prop:scatter}
For unit-norm embeddings $\hh_{k,i}\in\Sph^{d-1}$, define the average normalised
within-class scatter
$\sigma^2=\frac{1}{NK}\sum_{k=1}^{K}\sum_{i=1}^{N}\norm{\hh_{k,i}-\mmu_k}_2^{2}$.
Then
\begin{equation}
\sigma^2 = 1 + \Lscat .
\label{eq:equiv}
\end{equation}
Consequently, minimising $\Lscat$ is equivalent to minimising the within-class
scatter $\sigma^2$, and equivalently to maximising the mean intra-class kernel
alignment $\frac{1}{NK}\sum_k\sum_{i,j}\inner{\hh_{k,i}}{\hh_{k,j}}
=\frac{1}{K}\sum_k\norm{\mmu_k}_2^2$.
\end{proposition}

\begin{proof}
Fix a kernel $k$. Expanding the squared norm and using $\norm{\hh_{k,i}}_2^2=1$,
\begin{align}
\sum_{i=1}^{N}\norm{\hh_{k,i}-\mmu_k}_2^{2}
&=\sum_{i=1}^{N}\!\Big(\norm{\hh_{k,i}}_2^{2}
-2\inner{\hh_{k,i}}{\mmu_k}+\norm{\mmu_k}_2^{2}\Big)\nonumber\\
&=N-2N\norm{\mmu_k}_2^{2}+N\norm{\mmu_k}_2^{2}
=N\big(1-\norm{\mmu_k}_2^{2}\big),
\label{eq:proofstep}
\end{align}
where $\sum_i\inner{\hh_{k,i}}{\mmu_k}=N\inner{\mmu_k}{\mmu_k}=N\norm{\mmu_k}_2^2$
follows from the definition of $\mmu_k$. Averaging Eq.~\eqref{eq:proofstep} over
$k$ and dividing by $N$ yields
$\sigma^2=1-\frac{1}{K}\sum_k\norm{\mmu_k}_2^2=1+\Lscat$. The
alignment identity is immediate from
$\sum_{i,j}\inner{\hh_{k,i}}{\hh_{k,j}}=N^2\norm{\mmu_k}_2^2$.
\end{proof}

\begin{remark}[Connection to one-class CMKKM]
The one-class instance of centroid multiple-kernel $k$-means maximises
$\operatorname{Tr}(\bm{H}^{\!\top}\bm{K}_{\eta}\bm{H})$ with a single soft cluster
indicator $\bm{H}=N^{-1/2}\bm{1}$ and a convex kernel combination
$\bm{K}_{\eta}=\sum_k\eta_k\,\Phi_k\Phi_k^{\!\top}$, where
$\Phi_k=[\hh_{k,1},\dots,\hh_{k,N}]^{\!\top}$ realises a linear kernel on the
embeddings. A short calculation gives
$\operatorname{Tr}(\bm{H}^{\!\top}\bm{K}_{\eta}\bm{H})
=N\sum_k\eta_k\norm{\mmu_k}_2^2$, which under uniform weights
$\eta_k=1/K$ equals $-N\,\Lscat$. Hence $\Lscat$ is precisely the
(negated, scale-normalised) one-class CMKKM compactness objective, providing a
multiple-kernel-learning justification for the penalty.
\end{remark}

\subsection{Joint objective and collapse avoidance}
\label{subsec:joint}

The encoder $\{\WW_k\}_{k=1}^{K}$ is trained by minimising the combined loss
\begin{equation}
\Ltot=\Lcross+\lambda\,\Lscat,
\label{eq:total}
\end{equation}
where $\lambda\ge 0$ trades off cross-kernel consistency against intra-class
compactness; $\lambda=0$ recovers the purely contrastive backbone.

A natural concern is representational \emph{collapse}: the scatter term alone is
trivially minimised ($\Lscat=-1$) when all embeddings of a kernel coincide
($\norm{\mmu_k}_2=1$). This degenerate solution is, however, suppressed by the
contrastive term. In the alignment--uniformity decomposition of contrastive
learning, the numerator of Eq.~\eqref{eq:infonce} promotes \emph{alignment} of
positives while the log-sum-exp denominator promotes \emph{uniformity} of the
embedding distribution on the sphere, which is maximised only by a spread-out
configuration and is therefore directly opposed to collapse. The joint
objective~\eqref{eq:total} thus realises a controlled balance: $\Lscat$ pulls
each normal class towards its centroid, $\Lcross$ prevents the normal embeddings
from degenerating to a point and keeps different samples separable, and
$\lambda$ governs the equilibrium. The unit-sphere constraint additionally
bounds $\norm{\mmu_k}_2\le 1$, so the penalty cannot diverge.

\subsection{Anomaly scoring via linear one-class SVM}
\label{subsec:detection}

After training, we discard the contrastive heads' interaction and retain the
\emph{concatenated} representation
\begin{equation}
\zz_i=\big[\hh_{1,i}^{\!\top},\hh_{2,i}^{\!\top},\dots,
\hh_{K,i}^{\!\top}\big]^{\!\top}\in\R^{Kd},
\label{eq:concat}
\end{equation}
which fuses all $K$ kernel views into a single descriptor. A linear-kernel
OC-SVM is fitted on the normal embeddings $\{\zz_i : y_i=0\}$ by solving
\begin{equation}
\min_{\ww,\rho,\bm{\xi}}\ \
\tfrac{1}{2}\norm{\ww}_2^{2}-\rho
+\frac{1}{\nu N_{\mathrm{n}}}\sum_{i:\,y_i=0}\xi_i
\quad\text{s.t.}\quad
\inner{\ww}{\zz_i}\ge\rho-\xi_i,\ \ \xi_i\ge 0,
\label{eq:ocsvm}
\end{equation}
where $\nu\in(0,1]$ upper-bounds the fraction of normal training points allowed
to fall outside the decision region. The anomaly score is the negative signed
distance to the learned hyperplane,
\begin{equation}
a(\xx_i)=\rho-\inner{\ww}{\zz_i},
\label{eq:score}
\end{equation}
so that points lying on the far (origin) side of the hyperplane receive high
scores. We retain the \emph{linear} kernel deliberately: the contrastive and
scatter objectives jointly mould the normal class into a compact, nearly
linearly separable cone on the product of spheres $(\Sph^{d-1})^{K}$, for which a
linear separator is both sufficient and substantially cheaper than an RBF
OC-SVM.

\subsection{Overall algorithm and complexity}
\label{subsec:algorithm}

Algorithm~\ref{alg:scmk} summarises the complete training and scoring pipeline.
The dominant per-mini-batch cost is the cross-kernel contrastive term,
$O\!\big(K^2 B^2 d\big)$ for the $\binom{K}{2}$ pairwise affinity matrices, while
the scatter regularizer contributes only $O(KBd)$. Embedding extraction over the
full data set is $O(NKdD)$, and fitting the linear OC-SVM scales with the number
of support vectors. With a fixed small kernel bank ($K=5$), the method retains
the same asymptotic complexity as the contrastive backbone, so the accuracy gains
reported in the experiments come at no additional order-of-growth cost.

\begin{algorithm}[t]
\caption{SCMK for one-class anomaly detection}
\label{alg:scmk}
\begin{algorithmic}[1]
\Require normal set $\Xnorm$, full set $\mathcal{X}$, kernel ratios
$\{r_k\}_{k=1}^{K}$, embedding dim.\ $d$, weight $\lambda$, batch size $B$,
epochs $T$, OC-SVM parameter $\nu$
\Ensure anomaly scores $\{a(\xx_i)\}_{i=1}^{N}$
\State $\mathrm{med}\gets$ median pairwise distance on $\Xnorm$
\Comment{Eq.~\eqref{eq:median}}
\State build kernels $\kappa_k$ with $t_k=r_k\,\mathrm{med}$
\Comment{Eq.~\eqref{eq:gaussian}}
\State initialise projection heads $\{\WW_k\}_{k=1}^{K}$
\For{$\mathrm{epoch}=1$ \textbf{to} $T$}
  \For{each mini-batch of $B$ samples from $\Xnorm$}
    \State $\hh_{k,i}\gets \WW_k\xx_i/\norm{\WW_k\xx_i}_2$ for all $k,i$
    \Comment{Eq.~\eqref{eq:embedding}}
    \State $\Lcross\gets$ cross-kernel InfoNCE over kernel pairs
    \Comment{Eqs.~\eqref{eq:affinity}--\eqref{eq:lcross}}
    \State $\Lscat\gets -\tfrac{1}{K}\sum_k\norm{\tfrac{1}{B}\sum_i\hh_{k,i}}_2^2$
    \Comment{Eq.~\eqref{eq:lscat}}
    \State update $\{\WW_k\}$ by a gradient step on
    $\Lcross+\lambda\Lscat$
    \Comment{Eq.~\eqref{eq:total}}
  \EndFor
\EndFor
\State $\zz_i\gets[\hh_{1,i};\dots;\hh_{K,i}]$ for all $\xx_i\in\mathcal{X}$
\Comment{Eq.~\eqref{eq:concat}}
\State fit linear OC-SVM on $\{\zz_i:y_i=0\}$ with parameter $\nu$
\Comment{Eq.~\eqref{eq:ocsvm}}
\State \Return $a(\xx_i)=\rho-\inner{\ww}{\zz_i}$
\Comment{Eq.~\eqref{eq:score}}
\end{algorithmic}
\end{algorithm}

%% ======================================================================
%% Bibliography (placeholder entries for the Method section citations).
\begin{thebibliography}{00}

\bibitem{oord2018infonce}
A.~van den Oord, Y.~Li, O.~Vinyals,
\textit{Representation learning with contrastive predictive coding},
arXiv:1807.03748, 2018.

\bibitem{chen2020simclr}
T.~Chen, S.~Kornblith, M.~Norouzi, G.~Hinton,
\textit{A simple framework for contrastive learning of visual
representations}, in: ICML, 2020, pp.\ 1597--1607.

\bibitem{wang2020uniformity}
T.~Wang, P.~Isola,
\textit{Understanding contrastive representation learning through alignment
and uniformity on the hypersphere}, in: ICML, 2020, pp.\ 9929--9939.

\bibitem{gonen2011mkl}
M.~G\"onen, E.~Alpayd{\i}n,
\textit{Multiple kernel learning algorithms},
Journal of Machine Learning Research 12 (2011) 2211--2268.

\bibitem{huang2012mkkm}
H.-C.~Huang, Y.-Y.~Chuang, C.-S.~Chen,
\textit{Multiple kernel fuzzy clustering},
IEEE Transactions on Fuzzy Systems 20~(1) (2012) 120--134.

\bibitem{scholkopf2001ocsvm}
B.~Sch\"olkopf, J.~C.~Platt, J.~Shawe-Taylor, A.~J.~Smola, R.~C.~Williamson,
\textit{Estimating the support of a high-dimensional distribution},
Neural Computation 13~(7) (2001) 1443--1471.

\bibitem{ruff2018svdd}
L.~Ruff, R.~Vandermeulen, N.~G\"ornitz, et al.,
\textit{Deep one-class classification}, in: ICML, 2018, pp.\ 4393--4402.

\end{thebibliography}

\end{document}
